\section{Race 2 --- conditioning covariance on the label}
\label{sec:race2}

\aidraft{%
If timing exposure fails, perhaps the state helps estimate \emph{risk} rather than direction.
We build a long-only equal-risk-contribution portfolio over equity, ten-year Treasury total return,
gold, and cash, targeting 8\% volatility (scale-down only), and condition its covariance estimate on
the frozen state label (hard state, expanding per-state windows, 0.5 shrinkage to unconditional),
rebalancing monthly and on state flips.
The matched control is the identical allocator with conditioning removed; the reactive co-primary is
the same allocator on a plain EWMA covariance ($\lambda{=}0.97$).%
}

\aidraft{%
Conditioning adds essentially nothing over its matched twin --- $\feeAmatch$~bps/yr, CI
$\feeAmatchci$, inside all null bands (Table~\ref{tab:exp2}) --- and loses to the reactive
incumbent: $\text{fee}_B(\text{cond}-\text{EWMA}) = \feeBreact$~bps/yr, with the EWMA arm the best
performer outright (Table~\ref{tab:exp2arms}).
The mechanism the state is supposed to exploit is real but second-best: conditioning does stabilize
portfolio risk (vol-of-vol $\volofvolCond$ vs $\volofvolMatch$~pp for the matched twin), it is simply
dominated at the job by an estimator that re-estimates the whole covariance in weeks while the
expanding per-state estimate drags decades of stale mixture.
The frozen classification is null.%
}

\begin{table}[t]
\centering
\small
\caption{Race 2 --- conditional allocation: primary and controls. FKO fee, annualized bps/yr,
$\gamma{=}10$. Source: \texttt{results/allocation\_summary.csv}, \texttt{allocation\_controls.csv}.}
\label{tab:exp2}
\begin{tabular}{lr}
\toprule
Quantity & Value \\
\midrule
fee$_A$(cond $-$ B\_match), full window       & \feeAmatch~\feeAmatchci \\
fee$_B$(cond $-$ EWMA), from activation       & \feeBreact~\feeBreactci \\
fee$_A$ / fee$_B$ at $\gamma{=}1$             & $+1.2$ / $+3.2$ \\
placebo ($n{=}100$): fee$_A$ 5--95th          & $[-0.5,\,+5.8]$ \\
placebo ($n{=}100$): fee$_B$ 5--95th          & $[-32.7,\,-25.7]$ \\
episode-shuffle ($n{=}50$): fee$_A$ 5--95th   & $[-1.7,\,+4.9]$ \\
iid panels ($n{=}20$): fee$_A$ 5--95th        & $[-1.2,\,+0.9]$ \\
vol-of-vol, cond / B\_match (F2)              & \volofvolCond{} / \volofvolMatch{}~pp \\
turnover/yr, cond / B\_match / EWMA          & 0.11 / 0.06 / 0.88 \\
conditioning activation                       & 1994-03-29 \\
falsifiers F1 / F1b / F2 / F3                  & T / T / F / F \\
frozen classification                         & NULL \\
\bottomrule
\end{tabular}
\end{table}

\begin{table}[t]
\centering
\small
\caption{Race 2 --- arm summary over the scored window. The reactive EWMA arm is the best performer;
buy-and-hold equity is shown for exposure context. Source: \texttt{results/allocation\_arms.csv}.}
\label{tab:exp2arms}
\begin{tabular}{lrrrr}
\toprule
Arm & Ann.\ return & Ann.\ vol & Sharpe & MaxDD \\
\midrule
Conditional ERC              & 8.34\% & 6.88\%  & 0.826 & $-18.3\%$ \\
B\_match (uncond.\ expanding) & 8.33\% & 6.90\%  & 0.822 & $-18.3\%$ \\
B\_react (EWMA $\lambda{=}0.97$) & 8.28\% & 6.40\% & \textbf{0.878} & $-16.2\%$ \\
60/40                        & 9.44\% & 10.68\% & 0.634 & $-30.8\%$ \\
VT equity                    & 8.37\% & 10.04\% & 0.568 & $-26.4\%$ \\
B\&H equity                  & 12.30\% & 18.10\% & 0.533 & $-54.6\%$ \\
\bottomrule
\end{tabular}
\end{table}

